\section{Definitions}

Let $R$ be a commutative ring. Recall that a formal group law is defined to be a multivariate power series
$F \in R[\![X,Y]\!]$ satisfies following two conditions:
\begin{enumerate}
  \item $F(X,Y) = X + Y + \sum_{i, j \ge 1}^{\infty} c_{ij} X ^ i Y ^ j$.
  \item $F(F(X,Y), Z) = F(X, F(Y,Z))$.
\end{enumerate}

In this section, we fix some local notation for multivariate
power series variables using the tactic \lstinline|name_power_vars|.
For power series in two variables, we write \lstinline|X₀| and
\lstinline|X₁| for the two coordinate variables in
\lstinline|MvPowerSeries (Fin 2) R|.  For power
series in three variables, we write \lstinline|Y₀|, \lstinline|Y₁|, and
\lstinline|Y₂| for the three coordinate variables in
\lstinline|MvPowerSeries (Fin 3) R|.
In Lean, these notations can be introduced as follows:
\begin{lstlisting}
name_power_vars X₀, X₁ over R

name_power_vars Y₀, Y₁, Y₂ over R
\end{lstlisting}

A formal group laws is defined as
\href{https://leanprover-community.github.io/mathlib4_docs/Mathlib/RingTheory/FormalGroup/Basic.html#FormalGroup}{\faExternalLink}
\begin{lstlisting}
@[ext]
structure FormalGroup (R : Type*) [CommRing R] where
  /-- The underlying power series $F(X, Y)$ in two variables. -/
  toPowerSeries : MvPowerSeries (Fin 2) R
  /-- The constant coefficient of the formal group law is zero. -/
  zero_constantCoeff : toPowerSeries.constantCoeff = 0
  /-- The coefficient of $X$ in $F(X, Y)$ is 1. -/
  lin_coeff_X : toPowerSeries.coeff (single 0 1) = 1
  /-- The coefficient of $Y$ in $F(X, Y)$ is 1. -/
  lin_coeff_Y : toPowerSeries.coeff (single 1 1) = 1
  /-- Associativity condition: $F(F(X, Y), Z) = F(X, F(Y, Z))$. -/
  assoc : toPowerSeries.subst ![toPowerSeries.subst ![Y₀, Y₁], Y₂] =
    toPowerSeries.subst ![Y₀, toPowerSeries.subst ![Y₁, Y₂]] (S := R)
\end{lstlisting}

In mathlib, suppose that \lstinline|a b c| is of type \lstinline|α|, then
\lstinline|![a, b]| can be viewed as the function from \lstinline|Fin 2|
sending \lstinline|0| to \lstinline|a| and \lstinline|1| to \lstinline|b|.
Therefore, \lstinline|![Y₀, Y₁]| is a function from \lstinline|Fin 2| to
\lstinline|MvPowerSeries (Fin 2) R|, sending \lstinline|0| to
\lstinline|Y₀| and \lstinline|1| to \lstinline|Y₁|.

\begin{lstlisting}
/-- The natural inclusion from formal group law into formal power series. -/
instance FormalGroup.coeToPowerSeries :
  Coe (FormalGroup R) (MvPowerSeries (Fin 2) R) := ⟨toPowerSeries⟩

/-- Given a formal group `F`, `F.IsComm` is a proposition that $F(X,Y) = F(Y,X)$. -/
class FormalGroup.IsComm (F : FormalGroup R) : Prop where
  comm : F = (F : MvPowerSeries (Fin 2) R).subst ![X₁, X₀]
\end{lstlisting}

The instance \lstinline|FormalGroup.coeToPowerSeries| allows Lean to
coerce a formal group law to its underlying power series.  Thus, if
\lstinline|F : FormalGroup R|, then Lean can interpret \lstinline|F| as
a term of type \lstinline|MvPowerSeries (Fin 2) R| by using the field
\lstinline|F.toPowerSeries|.

The class \lstinline|FormalGroup.IsComm F| records the commutativity
condition for a formal group law.  It states that \lstinline|F| is equal
to the power series obtained by substituting \lstinline|X₁| for the first
variable and \lstinline|X₀| for the second variable.  In ordinary
mathematical notation, this says
\[
  F(X_0, X_1) = F(X_1, X_0).
\]

In the early stages of this project,
I define commutative formal group laws as a separate structure
\lstinline|CommFormalGroup| extending \lstinline|FormalGroup|.
There are some reasons for me to change this definition.

First, if commutative formal group laws were packaged
as a new structure, then a commutative formal group law and its underlying
formal group law would have different types.  As a result, many
definitions and lemmas stated for \lstinline|FormalGroup R| would require
additional coercions or separate restatements for
\lstinline|CommFormalGroup R|.

By contrast, the typeclass approach keeps the object itself as a term
\lstinline|F : FormalGroup R| and records commutativity as an additional
property \lstinline|[F.IsComm]|.  Thus all constructions and theorems for
general formal group laws can still be applied directly to \lstinline|F|,
while commutativity can be requested only in the places where it is
needed. Moreover, once an instance \lstinline|[F.IsComm]| is available, Lean can
find the commutativity assumption automatically by typeclass inference.

Recall that the additive inverse of a formal group law is a power series
$i(X)$ satisfying that $F(X, i(X)) = 0$. In Lean, we construct this power series
by induction as

\begin{lstlisting}
abbrev addInv_aux (F : FormalGroup R) : ℕ → R
  | 0 => 0
  | 1 => -1
  | n + 1 => - (coeff (n + 1) (F.toPowerSeries.subst
    ![X, (∑ i : Fin (n + 1), C (addInv_aux F i.1) * X ^ i.1)]))

def addInv_X : PowerSeries R := PowerSeries.mk (addInv_aux F ·)

def addInv (φ : MvPowerSeries σ R) : MvPowerSeries σ R := subst φ (addInv_X F)
\end{lstlisting}

Here \lstinline|addInv_aux F n| is the coefficient of degree \lstinline|n|
in the formal inverse of \lstinline|F|.  The first two coefficients are
defined to be \lstinline|0| and \lstinline|1|. For the higher coefficients,
we defines them recursively.  Suppose
the coefficients up to degree \lstinline|n| have already been constructed.
Then we have a truncated inverse series
\lstinline|∑ i : Fin (n + 1), C (addInv_aux F i.1) * X ^ i.1|
and substitutes it into the second variable of \lstinline|F|. And we define
the $n+1$-th coefficient of additive inverse to be the
minus $n+1$-th coefficient of this power series.

The function
\lstinline|PowerSeries.mk| constructs a one-variable formal
power series from a sequence of coefficients. The function
\lstinline|addInv φ| is a multivariate power series obtained by
substituting \lstinline|φ| into the power series \lstinline|addInv_X|.
Mathematically, it is the additive inverse of \lstinline|φ| in the sense
of formal group law \lstinline|F|.


We now use a formal group law \lstinline|F| to define a new group
structure on certain multivariate power series.  More precisely, we
consider power series satisfying \lstinline|PowerSeries.HasSubst|, which
means that their constant coefficient is nilpotent and hence they can be
substituted into another power series.  These power series will be called
the \emph{points} of the formal group law \lstinline|F|.

In Lean, this type is defined as follows:
\begin{lstlisting}
def Point (F : FormalGroup R) (σ : Type*) :=
  {f : MvPowerSeries σ R // PowerSeries.HasSubst f}
\end{lstlisting}

Here \lstinline|F.Point σ| is a subtype of \lstinline|MvPowerSeries σ R|.
An element \lstinline|x : F.Point σ|
consists of an underlying power series \lstinline|x.val| with a
proof \lstinline|x.prop| that it satisfies
\lstinline|PowerSeries.HasSubst|.

The formal group law \lstinline|F| then defines an addition operation on
these points.  If \lstinline|x y : F.Point σ|, their sum is obtained by
substituting \lstinline|x.val| and \lstinline|y.val| into the two
variables of \lstinline|F|.

In Lean this is written as:
\begin{lstlisting}
instance : Add (F.Point σ) where
  add x y := ⟨(F : MvPowerSeries (Fin 2) R).subst ![x.val, y.val], ...⟩
\end{lstlisting}

The expression
\lstinline|(F : MvPowerSeries (Fin 2) R).subst ![x.val, y.val]|
means that we regard \lstinline|F| as its underlying bivariate power
series and substitute \lstinline|x.val| for the first variable and
\lstinline|y.val| for the second variable.
The second component of the pair proves that the resulting power series
again satisfies \lstinline|PowerSeries.HasSubst|.

The zero element is given by the zero power series:
\begin{lstlisting}
instance : Zero (F.Point σ) where
  zero := ⟨0, PowerSeries.HasSubst.zero⟩
\end{lstlisting}

Thus \lstinline|F.Point σ| has a natural zero element and an addition
operation defined by the formal group law \lstinline|F|.  This addition
is not ordinary addition of power series; it is the operation obtained by
evaluating the formal group law on two points.

After defining addition and zero on \lstinline|F.Point σ|, Mathlib then
shows that these operations satisfy the usual algebraic laws.  The main
idea is that the formal group law axioms for \lstinline|F| induce the
corresponding group laws on the type of points \lstinline|F.Point σ|.

%%%%%%%%%%%%%%%%%%%%%%%%%%%%%%%%%%%%%%%%%%%

Now we discuss by steps to prove the group instance for
\lstinline|F.Point σ|.

The first step is to prove that the zero point is a left and right
identity for the addition defined by \lstinline|F|:
\begin{lstlisting}
instance : AddZeroClass (F.Point σ) where
  zero_add x := Subtype.ext (zero_add F x.prop)
  add_zero x := Subtype.ext (add_zero F x.prop)
\end{lstlisting}

Here \lstinline|zero_add| and \lstinline|add_zero| say that adding the
zero point on either side gives back the original point.  Since
\lstinline|F.Point σ| is a subtype, Lean proves equality of two points by
using \lstinline|Subtype.ext|, reducing the goal to an equality of their
underlying power series.

Next, associativity of the formal group law gives associativity of the
addition on points:
\begin{lstlisting}
instance : AddMonoid (F.Point σ) where
  nsmul := nsmulRec
  add_assoc x y z := Subtype.ext <| F.assoc' x.prop y.prop z.prop
\end{lstlisting}

The theorem \lstinline|F.assoc'| is the pointwise form of the associativity
condition of the formal group law.  It says that substituting three valid
points into the two sides of the associativity identity gives the same
power series.  Therefore the operation on \lstinline|F.Point σ| is
associative.  The field \lstinline|nsmulRec| tells Lean to use
the usual recursive definition of repeated addition by natural numbers.

If \lstinline|F| is commutative, then the addition on
\lstinline|F.Point σ| is also commutative:
\begin{lstlisting}
instance [F.IsComm] : AddCommMonoid (F.Point σ) where
  add_comm x y := Subtype.ext <| F.comm' x.prop y.prop
\end{lstlisting}

Here the assumption \lstinline|[F.IsComm]| means that the formal group law
itself satisfies \lstinline|F(X,Y) = F(Y,X)|.  The theorem
\lstinline|F.comm'| transfers this identity to points, proving that
\lstinline|x + y = y + x| for \lstinline|x y : F.Point σ|.

We could also records the semigroup structure separately:
\begin{lstlisting}
instance : AddSemigroup (F.Point σ) where
  add_assoc x y z := Subtype.ext <| F.assoc' x.prop y.prop z.prop
\end{lstlisting}

This contains only the associativity part of the structure.  Again, the
proof is obtained directly from the associativity of the formal group law
\lstinline|F|.

The next important step is the additive group structure:
\begin{lstlisting}
instance : Neg (F.Point σ) where
  neg f := ⟨F.addInv f.val, ...⟩

instance : AddGroup (F.Point σ) where
  nsmul := nsmulRec
  zsmul := zsmulRec
  neg_add_cancel x := ...
\end{lstlisting}

This instance says that every point has an additive inverse with respect
to the formal group law operation.  The proof of
\lstinline|neg_add_cancel| shows that \lstinline|-x + x = 0|.  The lemma
\lstinline|F.add_neg_cancel| is the corresponding inverse identity for
the formal group law.  The proof rewrites the expression using
associativity, the zero laws, and the inverse law until the desired
identity follows.  The fields \lstinline|nsmul := nsmulRec| and
\lstinline|zsmul := zsmulRec| specify the usual recursive definitions of
multiplication by natural numbers and integers.

Finally, if \lstinline|F| is commutative, the additive group is upgraded
to an additive commutative group:
\begin{lstlisting}
instance [F.IsComm] : AddCommGroup (F.Point σ) where
  add_comm x y := Subtype.ext <| F.comm' x.prop y.prop
\end{lstlisting}

Thus the formal group law \lstinline|F| gives a natural group structure
on \lstinline|F.Point σ|.  Without assuming commutativity, the points form
an additive group.  If the additional typeclass assumption
\lstinline|[F.IsComm]| is available, then this group is an additive
commutative group.



%%%%%%%%%%%%%%%%%%%%%%%%%%%%%%%%%%%%%%%%%%%%%%%%%%%%%%%%%%%%%%%%%%%%%%%%%%%%

% Given a formal group law, we can define a group structure on
% \lstinline{MvPowerSeries σ R}. To be precise, let \lstinline{(F : FormalGroup R) (f₀ f₁ : MvPowerSeries σ R)}.
% The sum of \lstinline{f₀} and \lstinline{f₁} is defined as \lstinline{subst ![f₀, f₁] F.toFun}.
% In Lean, we can introduce the simpler notation

% \begin{lstlisting}
% abbrev add (F : FormalGroup R) (f g : MvPowerSeries σ R) : MvPowerSeries σ R := subst ![f, g] F.toFun
% \end{lstlisting}

% In the environment of \lstinline{namespace FormalGroup}, \lstinline{+[F]} denotes addition with respect to \lstinline{F : FormalGroup R}. More specifically, given two multivariate power series $f, g \in R[\![X_I]\!]$, where $I$ is an index set, and assuming the constant coefficients of $f$ and $g$ are zero,
% \lstinline{f +[F] g} means $F(f,g)$.

% With this notation, we can rewrite the associativity condition and the commutativity
% condition in the \lstinline{namespace FormalGroup} as follows:
% \begin{lstlisting}
% /-- The addition under the sense of formal group `F` is associative. -/
% theorem add_assoc {F : FormalGroup R} {Z₀ Z₁ Z₂ : MvPowerSeries σ R} (hZ₀ : constantCoeff Z₀ = 0) (hZ₁ : constantCoeff Z₁ = 0) (hZ₂ : constantCoeff Z₂ = 0):
% Z₀ +[F] Z₁ +[F] Z₂ = Z₀ +[F] (Z₁ +[F] Z₂) := ...
% \end{lstlisting}

% \begin{lstlisting}
% theorem add_comm {F : FormalGroup R} [hF : F.IsComm] {Z₀ Z₁ : MvPowerSeries σ R}
% (hZ₀ : constantCoeff Z₀ = 0) (hZ₁ : constantCoeff Z₁ = 0):
%   Z₀ +[F] Z₁ = Z₁ +[F] Z₀ := ...
% \end{lstlisting}

% To show that the formal addition defined by a formal group law $F$ forms an
% additive group, we should verify that there is a unit for this addition, namely
% $F(X,0)=X$ and $F(0,Y)=Y$. Furthermore, we should prove that there exists a unique additive
% inverse element. \par

% The existence of an additive unit in the sense of $F$ is formalized by the following two theorems.
% \begin{lstlisting}
% theorem zero_left : subst ![0, PowerSeries.X] F.toFun = PowerSeries.X (R := R) := ...

% theorem zero_right : subst ![PowerSeries.X, 0] F.toFun = PowerSeries.X (R := R) := ...
% \end{lstlisting}

% Using the same notation for formal group addition, we can generalize the above theorems to
% the following two theorems.

% \begin{lstlisting}
% theorem zero_add (h : constantCoeff f = 0) : 0 +[F] f = f :=

% theorem add_zero (h : constantCoeff f = 0) : f +[F] 0 = f :=
% \end{lstlisting}

% Here, the name \lstinline{zero_add} means that adding zero to a multivariate power series
% \lstinline{f} gives \lstinline{f} itself, which is a standard naming convention in Mathlib.
% To distinguish this theorem from other theorems with the same name, we place it in the
% \lstinline{namespace FormalGroup}.

% Now we discuss the construction of the additive inverse of a formal group law.

% \begin{lstlisting}
% abbrev addInv_aux (F : FormalGroup R) : ℕ → R
%   | 0 => 0
%   | 1 => -1
%   | n + 1 => - (coeff (n + 1) (F.toFun.subst
%     ![X, (∑ i : Fin (n + 1), C (addInv_aux F i.1) * X ^ i.1)]))

% def addInv_X : PowerSeries R := .mk (addInv_aux F ·)

% def addInv (φ : MvPowerSeries σ R) : MvPowerSeries σ R := subst φ (addInv_X F)
% \end{lstlisting}

% The API \lstinline|addInv_X| is the additive inverse of \lstinline|X| with respect to the formal group
% law \lstinline|F|. For an arbitrary multivariate power series, \lstinline|addInv φ| is the additive
% inverse of \lstinline|φ|. Via direct calculation and induction, we have the theorem

% \begin{lstlisting}
% theorem subst_addInv_eq_zero : F.toFun.subst ![X, (addInv_X F)] = 0 := by
% \end{lstlisting}

\subsection{Formal groups constructions}

The two first examples of formal group laws are the additive formal group law
$\GG_a$ and the
multiplicative formal group law $\GG_m$. We formalize them using
our definition of \lstinline{FormalGroup R} and show that
they are commutative.
\begin{lstlisting}
def 𝔾ₐ : FormalGroup R where
  toPowerSeries := X₀ + X₁

instance : (𝔾ₐ (R := R)).IsComm where ...

def 𝔾ₘ : FormalGroup R where
  toPowerSeries := X₀ + X₁ + X₀ * X₁

instance : (𝔾ₘ (R := R)).IsComm where ...
\end{lstlisting}

% \subsection{Functions Between Formal Groups}

% According to \ref{defn: pushforward of coeff ring},

Another way to construct new formal group laws is to change the
coefficient ring of an existing one.
Recall that given two commutative rings $R$ and $R'$, a ring
homomorphism $f : R \to R'$, and a formal group law
\lstinline|F| over $R$, the coefficientwise image of \lstinline|F| under
$f$ is again a formal group law over $R'$.  This new formal group
law is usually denoted by $f_* F$.

The construction of this pushforward on the coefficient ring is defined as
\href{https://github.com/WenrongZou/FormalGroupLaws/blob/main/FormalGroupLaws/Basic.lean}{\faExternalLink}

\begin{lstlisting}
def map {R : Type*} [CommRing R] {R' : Type*} [CommRing R']
    (f : R →+* R') (F : FormalGroup R) : FormalGroup R' where
  toPowerSeries := F.toPowerSeries.map f
\end{lstlisting}

There is also a converse kind of construction.  Sometimes a formal group
law is originally defined over a larger ring $R$, but all of its
coefficients actually lie in a smaller subring $T$.  In this situation,
the formal group law can be regarded as being defined over $T$.

In other words, given a formal group law \lstinline|F| over a ring $R$,
if every coefficient of \lstinline|F| belongs to a subring $T$, then
\lstinline|F| descends to a formal group law over $T$.  This is useful
in our formalization of the functional equation lemma.

The construction of this is defined as:

\begin{lstlisting}
def FormalGroup.toSubring (F : FormalGroup R) (T : Subring R)
    (hF : ∀ n, F.toFun n ∈ T) : FormalGroup T where
  toPowerSeries := F.toPowerSeries.toSubring _ hF
\end{lstlisting}

The function \lstinline|MvPowerSeries.toSubring|
\href{https://leanprover-community.github.io/mathlib4_docs/Mathlib/RingTheory/MvPowerSeries/Basic.html#MvPowerSeries.toSubring}{\faExternalLink}
is used to regard a power
series whose coefficients lie in a subring as a power series over that
subring.  In this way, it changes the coefficient ring from the
ring $R$ to a subring $T$, provided that all coefficients of the power
series actually belong to $T$.

Moreover, the constructions above are compatible with commutativity.
More precisely, we prove that if the original formal group law
\lstinline|F| is commutative, then the formal group laws obtained by
changing the coefficient ring are also commutative.  In Lean, this gives
instances showing that, under the assumption \lstinline|F.IsComm|, both
\lstinline|(F.map f).IsComm| and \lstinline|(F.toSubring T hF).IsComm|
hold.

\subsection{Homomorphism and isomorphism}

Let $R$ be a commutative ring and $F,G$ be two formal group laws over $R$.
Recall that a formal group homomorphism $f$ from $F$ to $G$ is defined to be
a power series with constant coefficient such that
\[f (F(X,Y)) = G(f(X), f(Y)). \]

It is defined in Lean as

\begin{lstlisting}
@[ext]
structure FormalGroupHom (F G : FormalGroup R) where
  /-- The underlying power series of a formal group homomorphism. -/
  toPowerSeries : PowerSeries R
  /-- Constant coefficient of underlying power series is zero. -/
  zero_constantCoeff : toPowerSeries.constantCoeff = 0
  /-- The homomorphism condition: $f(F(X,Y))=G(f(X),f(Y))$. -/
  hom : toPowerSeries.subst F =
    G.toPowerSeries.subst (toPowerSeries.toMvPowerSeries ·)
\end{lstlisting}

Here \lstinline|toPowerSeries| is a power series, while
\lstinline|G.toPowerSeries| is a bivariate power series.  Therefore,
in the expression
\lstinline|G.toPowerSeries.subst (toPowerSeries.toMvPowerSeries ·)|, we
need to regard the one-variable power series \lstinline|toPowerSeries| as
a multivariable power series in each variable of \lstinline|G|.
This is done using \lstinline|PowerSeries.toMvPowerSeries|, whose
definition is

\begin{lstlisting}
def PowerSeries.toMvPowerSeries (i : σ) : PowerSeries R →ₐ[R] MvPowerSeries σ R :=
  MvPowerSeries.rename (fun _ => i)
\end{lstlisting}

The map \lstinline|PowerSeries.toMvPowerSeries|
is defined to be a special case of
\lstinline|MvPowerSeries.rename|
\href{https://leanprover-community.github.io/mathlib4_docs/Mathlib/RingTheory/MvPowerSeries/Rename.html#MvPowerSeries.rename}{\faExternalLink}.
This operation renames the variables of a multivariable power series
according to a map between indexing types.  In this case, the map
\lstinline|fun _ => i|
sends the variable of a power series to the variable
\lstinline|i| in \lstinline|MvPowerSeries σ R|.


Thus, in the homomorphism condition
\begin{lstlisting}
hom : toPowerSeries.subst F =
  G.toPowerSeries.subst (toPowerSeries.toMvPowerSeries ·)
\end{lstlisting}

The left-hand side represents $f(F(X,Y))$, while the right-hand side
represents $G(f(X), f(Y))$. And the usage of the tag \lstinline|@[ext]|
can be found in Section~\ref{sec: tag_ext_simp}.

A formal group isomorphism is a formal group homomorphism which has an
inverse formal group homomorphism.  In Lean, this is packaged as the
following structure:
\begin{lstlisting}
@[ext]
structure FormalGroupIso (F G : FormalGroup R) where
  /-- The underlying formal group homomorphism of a formal group isomorphism. -/
  toHom : FormalGroupHom F G
  /-- The inverse homomorphism of underlying formal group homomorphism. -/
  invHom : FormalGroupHom G F
  /-- `toHom ∘ invHom = id`. -/
  left_inv : toHom.toPowerSeries.subst ∘ (PowerSeries.subst invHom.toPowerSeries) = id
  /-- `invHom ∘ toHom = id`. -/
  right_inv : invHom.toPowerSeries.subst ∘ (PowerSeries.subst toHom.toPowerSeries) = id
\end{lstlisting}

Thus \lstinline|FormalGroupIso F G| bundles together two homomorphisms:
\lstinline|toHom : FormalGroupHom F G| and
\lstinline|invHom : FormalGroupHom G F|.
The fields \lstinline|left_inv| and \lstinline|right_inv| state that
these two homomorphisms are inverse to each other under composition of
power series.  More precisely, \lstinline|left_inv| says that composing
\lstinline|invHom| and then \lstinline|toHom| gives the identity, while
\lstinline|right_inv| says that composing \lstinline|toHom| and then
\lstinline|invHom| gives the identity.

Therefore, a term of type \lstinline|FormalGroupIso F G| is not merely a
map from \lstinline|F| to \lstinline|G|.  It contains a homomorphism, an
inverse homomorphism, and proofs that the two compositions are the
identity.

% An formal group isomorphism from $F$ to $G$ is defined as
% \begin{lstlisting}
% @[ext]
% structure FormalGroupIso (F G : FormalGroup R) where
%   /-- The underlying formal group homomorphism of a formal group isomorphism. -/
%   toHom : FormalGroupHom F G
%   /-- The inverse homomorphism of underlying formal group homomorphism. -/
%   invHom : FormalGroupHom G F
%   /-- `toHom ∘ invHom = id`. -/
%   left_inv : toHom.toPowerSeries.subst ∘ (PowerSeries.subst invHom.toPowerSeries) = id
%   /-- `invHom ∘ toHom = id`. -/
%   right_inv : invHom.toPowerSeries.subst ∘ (PowerSeries.subst toHom.toPowerSeries) = id
% \end{lstlisting}

Recall that a homomorphism $f$ from $F$ to $G$ is an isomorphism if and only
if the coefficient of $X$ in $f$ is a unit in $R$. This result is formalized
as:

\begin{lstlisting}
theorem isIso_of_isUnit_coeff_one {F G : FormalGroup R} (f : FormalGroupHom F G)
    (h : IsUnit (f.toPowerSeries.coeff 1)) :
    ∃ (g : FormalGroupIso F G), g.toHom = f := ...

theorem isUnit_coeff_one_of_isIso {F G : FormalGroup R} (f : FormalGroupIso F G) :
    IsUnit (f.toHom.toPowerSeries.coeff 1) := ...
\end{lstlisting}

There are two ways to express an element \lstinline|a| of a ring \lstinline|R| to
be invertible. The first one is \lstinline|IsUnit a|, while the second one
is \lstinline|Invertible a|. The type \lstinline|IsUnit a| is defined as a \lstinline|Prop|.
On the other hand, \lstinline|Invertible a| is a typeclass (residing in \lstinline|Type|)
that explicitly packages the element $a$ with its concrete inverse.

At the end of this section, we discuss the formalization of
substitution inverse of a power series. Let $R$ be a commutative ring.
Recall that given a power series $f$ over $R$, suppose that
the coefficient of $X$ in $f$ is invertible in $R$, then there is a
power series $g$ over $R$ such that $f(g(X)) = g(f(X)) = X$. This
$g$ is called substitution inverse of $f$.

In mathlib, it is defined as \lstinline|substInv|
\href{https://leanprover-community.github.io/mathlib4_docs/Mathlib/RingTheory/PowerSeries/Substitution.html#PowerSeries.substInv}{\faExternalLink}
and
\lstinline|substInvOfIsUnit|
\href{https://leanprover-community.github.io/mathlib4_docs/Mathlib/RingTheory/PowerSeries/Substitution.html#PowerSeries.substInvOfIsUnit}{\faExternalLink}
inside the \lstinline|namespace PowerSeries|.
\lstinline|substInvOfIsUnit| is a variant of \lstinline|substInv| using
assumption \lstinline|IsUnit| rather than \lstinline|Invertible|.


% \begin{lstlisting}
% theorem isUnit_of_invertible [Monoid α] (a : α) [Invertible a] : IsUnit a :=
% \end{lstlisting}

% The theorem \lstinline|isUnit_of_invertible| in Mathlib gives one direction of the
% equivalence between these two notions.

% In the proof of these two theorems, we need the result that, given a power series $f(X)$
% over a ring $R$, if the coefficient of $X$ in $f(X)$ is invertible and the constant coefficient
% is zero, then there exists a power series $g(X)$ over $R$ such that $f(g(X)) = g(f(X)) = X$.

% \begin{lstlisting}
% variable (P : PowerSeries R) (hP : P.constantCoeff = 0) [Invertible (coeff 1 P)]

% noncomputable
% def substInvFun : ℕ → R
%   | 0 => 0
%   | 1 => 1 / (coeff _ 1 P)
%   | n + 1 => - 1 / (coeff 1 P) *
%       (coeff (n + 1) (P.subst (∑ i : Fin (n + 1), C (substInvFun i.1) * X ^ i.1)))
% \end{lstlisting}

% From the construction and computation, we obtain the following theorem:

% \begin{lstlisting}
% lemma substInv_subst : P.substInv.subst P = X :=
% \end{lstlisting}

% Also, we can prove that this substitution inverse is unique.